% ============================================================
% ABSTRACT
% File: frontmatter/abstract.tex
% ============================================================

Industrial monitoring increasingly relies on data from cameras, acoustic signals, and
machine sensors. Although multimodal learning can leverage the strengths of each data
type, adding sources does not necessarily improve prediction accuracy or reliability,
particularly when data quality or relevance varies.

A critical review of 94 full-text works found that 46 directly tested the value of
adding an input, while 21 tested corrupted inputs and 11 tested missing inputs at
inference. This thesis investigates when visual information improves prediction beyond
machine sensors alone and whether fusion remains reliable when either stream
deteriorates.

An audit of Intel Welding found acquisition sessions shared across development and
assessment partitions, precluding independent welding validation. Four main studies
evaluated engineered sensor features, pretrained visual features, and patch-based
temporal Transformers using dataset-specific group splits. Fusion methods included
equal averaging, logistic stacking, and weighting by input quality.

Across 20 PHM2026 gear runs, combining three sensor groups to predict a photographic
proxy for gear damage yielded a mean absolute error of 0.824 percentage points,
compared with 0.804 for the training mean. The proxy lacked independent physical
validation. On REASSEMBLE, stacking video and sensor predictions increased pooled
failure average precision (AP) from 0.429 for sensors alone to 0.447 across 4,530
completed actions from 148 recordings. Its advantage over equal averaging (0.446)
remained unresolved. Severe sensor noise reduced stacking AP to 0.222 and impaired
calibration.

Sensor-only logistic regression outperformed stacking across all six MSM milling
protocols, with abnormal-cutting AP of 0.829 versus 0.757 on the primary cohort of
1,895 intervals. On ViFailback, evaluation across 4,382 episodes from 99 tasks held out
entire tasks. Class-averaged AP was 0.649 for sensors versus 0.621--0.629 for stacking
across three seeds. On 151 official HANDOVER test trials, sensors also ranked failures
better than equal averaging.

Visual information improved failure ranking in robotic assembly, but that gain did not
extend to milling or manipulation under the tested protocols. Adaptive weighting
offered no consistent protection against input corruption. These findings concern
retrospective condition estimation and completed-event classification. Early warning
and cross-factory deployment require separate validation.
